% ====================================================================
% PRÉAMBULE — moteur XeLaTeX (tectonic)
% ====================================================================
\usepackage{fontspec}
% police principale = Latin Modern par défaut de fontspec (pas besoin de la nommer)
\setsansfont{DejaVu Sans}
\setmonofont{DejaVu Sans Mono}[Scale=0.85]

\usepackage{polyglossia}
\setdefaultlanguage{french}
% guillemets français (non fournis par polyglossia)
\providecommand{\og}{«\,}
\providecommand{\fg}{\,»}

\usepackage{geometry}
\geometry{a4paper,margin=2.3cm,headheight=15pt}

\usepackage{amsmath,amssymb}
\usepackage{graphicx}
\usepackage{booktabs}
\usepackage{array}
\usepackage{multirow}
\usepackage{enumitem}
\usepackage{caption}
\captionsetup{font=small,labelfont=bf}

\usepackage[table]{xcolor}
% --- palette ---
\definecolor{bleu}{HTML}{2C3E72}
\definecolor{bleuclair}{HTML}{4C72B0}
\definecolor{vert}{HTML}{2E7D32}
\definecolor{vertclair}{HTML}{E8F5E9}
\definecolor{rouge}{HTML}{C0392B}
\definecolor{rougeclair}{HTML}{FDEDEC}
\definecolor{orange}{HTML}{E67E22}
\definecolor{orangeclair}{HTML}{FEF5E7}
\definecolor{violet}{HTML}{6C3483}
\definecolor{violetclair}{HTML}{F4ECF7}
\definecolor{gristresclair}{HTML}{F5F5F5}
\definecolor{griscode}{HTML}{2B2B2B}

% --- code Python ---
\usepackage{listings}
\definecolor{cmt}{HTML}{6A9955}
\definecolor{kw}{HTML}{569CD6}
\definecolor{str}{HTML}{CE9178}
\definecolor{num}{HTML}{B5CEA8}
\lstdefinestyle{py}{
  language=Python,
  basicstyle=\ttfamily\footnotesize,
  keywordstyle=\color{kw}\bfseries,
  commentstyle=\color{cmt}\itshape,
  stringstyle=\color{str},
  numberstyle=\tiny\color{gray},
  numbers=left,numbersep=7pt,
  backgroundcolor=\color{gristresclair},
  frame=leftline,framerule=2pt,rulecolor=\color{bleuclair},
  breaklines=true,breakatwhitespace=true,
  showstringspaces=false,
  tabsize=4,
  upquote=true,
  columns=fullflexible,keepspaces=true,
  literate=%
    {é}{{\'e}}1 {è}{{\`e}}1 {ê}{{\^e}}1 {ë}{{\¨e}}1
    {à}{{\`a}}1 {â}{{\^a}}1 {î}{{\^i}}1 {ï}{{\¨i}}1
    {ô}{{\^o}}1 {ö}{{\¨o}}1 {û}{{\^u}}1 {ù}{{\`u}}1 {ç}{{\c c}}1
    {œ}{{\oe}}1 {Œ}{{\OE}}1 {’}{{'}}1
    {É}{{\'E}}1 {È}{{\`E}}1 {Ê}{{\^E}}1 {À}{{\`A}}1
    {≤}{{$\leq$}}1 {≥}{{$\geq$}}1 {²}{{$^2$}}1 {→}{{$\rightarrow$}}1
    {≈}{{$\approx$}}1 {μ}{{$\mu$}}1 {≠}{{$\neq$}}1 {∈}{{$\in$}}1
    {├}{{|}}1 {└}{{|}}1 {─}{{-}}1 {│}{{|}}1 {’}{{'}}1 {«}{{<<}}1 {»}{{>>}}1
}
\lstset{style=py}

% --- boîtes colorées (tcolorbox) ---
\usepackage[most]{tcolorbox}
\tcbset{
  boxedstyle/.style={
    enhanced,breakable,boxrule=0pt,
    leftrule=3pt,arc=2pt,left=8pt,right=8pt,top=6pt,bottom=6pt,
    fonttitle=\bfseries,coltitle=white,
  }
}
\newtcolorbox{Definition}[1][]{boxedstyle,colback=violetclair,colframe=violet,
  title={\faIcon~Définition #1}}
\newtcolorbox{Exemple}[1][]{boxedstyle,colback=vertclair,colframe=vert,
  title={Exemple — #1}}
\newtcolorbox{Calcul}[1][]{boxedstyle,colback=gristresclair,colframe=bleuclair,
  title={Calcul détaillé #1}}
\newtcolorbox{Astuce}[1][]{boxedstyle,colback=orangeclair,colframe=orange,
  title={Astuce / À retenir #1}}
\newtcolorbox{Attention}[1][]{boxedstyle,colback=rougeclair,colframe=rouge,
  title={Attention #1}}
\newtcolorbox{Resume}[1][]{boxedstyle,colback=vertclair,colframe=vert,
  title={En résumé #1}}

% remplace \faIcon si fontawesome absent
\providecommand{\faIcon}[1]{}

% --- titres & en-têtes ---
\usepackage{titlesec}
\titleformat{\chapter}[display]
  {\normalfont\huge\bfseries\color{bleu}}
  {\filright\Large\color{bleuclair}Chapitre~\thechapter}{8pt}{\Huge}
\titlespacing*{\chapter}{0pt}{-20pt}{20pt}

\usepackage{fancyhdr}
\pagestyle{fancy}
\fancyhf{}
\fancyhead[L]{\small\color{bleu}\leftmark}
\fancyhead[R]{\small\thepage}
\renewcommand{\headrulewidth}{0.6pt}

\usepackage{hyperref}
\hypersetup{
  colorlinks=true,linkcolor=bleu,urlcolor=bleuclair,
  citecolor=vert,pdftitle={Cours de Data Mining — Explications complètes},
  pdfauthor={Master 2 IG/ID — ENEAM/UAC}
}

% --- commandes pratiques ---
\newcommand{\gini}{\textit{Gini}}
\newcommand{\R}{\mathbb{R}}
